\section{Arithmetic reduction and exact cosets}\label{sec:arithmetic}

The proof will show that certain finite quotients of a special unitary
rational group are trivial.  We first explain why these quotients suffice,
and why every one of their cosets admits the local approximations used
later.  Throughout, a global function field has its full finite field of
constants, and all quadratic field extensions are separable.  We use
normalized absolute values at its places.  A product indexed by the empty
set is the trivial group.

\subsection{The remaining unitary groups}

We use three established results to isolate the case that requires proof.
For a simply connected absolutely almost simple isotropic group over a
global field, the global Kneser--Tits theorem and Tits's simplicity theorem
give perfection and simplicity modulo the center of the rational group
\cite[Theorem~8.1]{Gille}\cite{Tits}.  A noncentral normal subgroup
therefore has central, hence abelian, quotient and must be the whole
rational group.  In this case the algebraic group is isotropic
at every completion, so there are no anisotropic places.  The
Margulis--Platonov theorem is also known for every inner form of type~$A$
over every global field
\cite[Theorem~6.4]{RapinchukSegevSeitz2002}; see also
\cite[Proposition~1]{Rapinchuk2006}.  Finally, Harder's theorem implies
that an anisotropic absolutely almost simple group over a global function
field has inner or outer type~$A$
\cite{Harder}\cite[p.~914]{PrasadTriality}.
Consequently it remains to consider an anisotropic outer form of type~$A$.

We write such a group as
\begin{equation}\label{arith:unitary-notation}
 S=\SU(D,*),\qquad U=\U(D,*),\qquad \deg_L D=n\geq 3,
\end{equation}
where $D$ is a central simple algebra over a quadratic field extension
$L/k$ and $*$ is an involution inducing the nontrivial automorphism of
$L/k$.  The algebra $D$ need not be a division algebra.  Degree two gives
an inner form of type~$A_1$, already covered by the preceding result.
Write $\bar a$ for the conjugate of $a\in L$ and
\[
 L^1=\{a\in L^\times:a\bar a=1\}.
\]
On $U$, the notation $\det$ will always mean the reduced norm
$\Nrd_{D/L}$.  It takes values in the norm-one torus
$\Res^1_{L/k}\bbG_m$, whose group of rational points is $L^1$.
We retain the notation
\[
 A=\{v:\rank_{k_v}S=0\},\qquad
 H_A=\prod_{v\in A}S(k_v),\qquad
 \delta_A:S(k)\longrightarrow H_A.
\]
The set $A$ is finite: outside finitely many places a reductive model has
quasi-split fibers and positive local rank.  Each factor of $H_A$ is
compact.

\begin{proposition}\label{arith:unitary-reduction}
For the groups in \eqref{arith:unitary-notation}, the following statements
hold in every positive characteristic.
\begin{enumerate}
\item The determinant maps $U(k)\to L^1$ and their local analogues are
surjective.
\item The groups $S$ and $U$ satisfy weak approximation.
\item Every place in $A$ splits in $L/k$.  At such a place there is a
central division algebra $\mathcal D_v$ of degree $n$ over $k_v$ for which
\[
 U(k_v)\simeq\mathcal D_v^\times,
 \qquad S(k_v)\simeq\SL_1(\mathcal D_v).
\]
Under this identification the determinant is $\Nrd_{\mathcal D_v/k_v}$.
\end{enumerate}
\end{proposition}

\begin{proof}
The exact sequence
\[
 1\longrightarrow S\longrightarrow U
 \xrightarrow{\det}\Res^1_{L/k}\bbG_m\longrightarrow1
\]
identifies each determinant fiber with an $S$-torsor.  Harder's global
vanishing theorem and the local vanishing theorem of Bruhat--Tits give
$H^1(k,S)=1$ and $H^1(k_v,S)=1$.  These statements include characteristic
two; a convenient uniform reference is
\cite[Theorem~5.1.1(i)]{Conrad}, which recalls
\cite[Satz~A]{Harder} and \cite[Theorem~4.7(ii)]{BruhatTitsIII}.
They prove (i).

Strong approximation for simply connected almost simple groups over
global function fields holds away from a place where the group is
isotropic \cite[Theorem~A]{PrasadSA}.  There are infinitely many split
places, by Chebotarev applied to a finite splitting field.  To approximate
at a given finite set of places, omit a split place outside that set.
This proves weak approximation for $S$.
For $U$, fix $\theta\in L\setminus k$.  The Cayley map
\begin{equation}\label{arith:cayley}
 x\longmapsto (x+\theta)(x+\bar\theta)^{-1},\qquad x^*=x,
\end{equation}
is a birational map from the affine space of Hermitian elements of $D$
to $U$.  Its inverse, on the open set where $u-1$ is invertible, is
$(u-1)^{-1}(\theta-u\bar\theta)$.  The difference
$\theta-\bar\theta$ is nonzero also in characteristic two.  Weak
approximation for this affine space, and density of the Cayley open set
in each local group, prove (ii).

Suppose that $v$ does not split in $L$.  A unitary involution implies
that the Brauer corestriction of $D$ is zero.  For a quadratic extension
of nonarchimedean local fields, corestriction preserves the local Brauer
invariant.  Thus $D\otimes_L L_v$ is a matrix algebra.  The group $S_{k_v}$
is consequently the special unitary group of an $n$-dimensional Hermitian
form.  The local classification of Hermitian forms over a separable
quadratic extension gives anisotropic dimension at most two, in all
characteristics.  Since $n\geq3$, this group is isotropic.  These standard
facts about unitary involutions and local Hermitian forms may be found in
\cite{KMRT,Reiner}; for the characteristic-two classification see also
\cite[Theorem~3.3 and Corollary~3.4]{ElomaryTignol}.

At a split place, the two simple components of $D\otimes_k k_v$ are
interchanged by the involution.  Choosing one component $B_v$ identifies
$U(k_v)$ with $B_v^\times$ and $S(k_v)$ with $\SL_1(B_v)$.  If
$B_v\simeq M_r(\mathcal D_v)$, this last group has rank $r-1$.
It is anisotropic exactly when $r=1$, proving (iii).
\end{proof}

\subsection{Finite index and the power quotient}

We will use the local Kneser--Tits theorem in the following form.  If a
simply connected absolutely almost simple group over a nonarchimedean
local field is isotropic, its group of local points is perfect and is
simple modulo its finite center.  In particular, a noncentral normal
subgroup is the whole group.  The perfection clause matters: simplicity
modulo the center first gives a central quotient, and perfection then
makes that quotient trivial.  This local result applies in equal
characteristic, including characteristics two and three
\cite{Tits,BruhatTitsIII,Gille}.

\begin{lemma}\label{arith:finite-index}
Every noncentral abstract normal subgroup of $S(k)$ has finite index.
\end{lemma}

\begin{proof}
This is the finite-index reduction in Margulis's normal subgroup theorem
\cite[Theorem~2.4.6 and Corollary~2.4.9]{Margulis1979}\cite{Margulis}.
We recall the passage from an arithmetic lattice to all rational points.
Let $N\triangleleft S(k)$ contain a noncentral element $z$.  Choose a
finite set of places $T$ containing $A$, the places where $z$ is not
integral, and two further places at which $S$ is split.  Fix compact open
subgroups $K_v\subset S(k_v)$ for $v\notin T$, using an integral model
outside finitely many places, and arrange that $z\in K_v$ there.  Put
\[
 J_T=\prod_{v\notin T}'S(k_v),\qquad
 K^T=\prod_{v\notin T}K_v,\qquad
 \Gamma=S(k)\cap K^T.
\]
Here $S(k)\cap K^T$ denotes the inverse image under the diagonal map to
$J_T$, and the prime denotes the restricted product with respect to the
$K_v$.

Reduction theory makes $\Gamma$ an arithmetic lattice in
$\prod_{v\in T}S(k_v)$; compact factors may be projected away.  Strong
approximation, omitting any one noncompact factor, gives irreducibility.
The two added split places have total rank at least four, since
$n\geq3$.  The arithmetic normal subgroup theorem therefore applies.
The subgroup $\Gamma$ is Zariski dense, so a finite normal subgroup of
$\Gamma$ is central.  Since $z\in N\cap\Gamma$ is noncentral, the theorem
gives
\begin{equation}\label{arith:lattice-index}
 [\Gamma:N\cap\Gamma]<\infty.
\end{equation}
The lattice and strong-approximation assertions used here are valid over
function fields \cite{HarderReduction,PrasadSA,Margulis}.

Let $C$ be the closure of $N$ in $J_T$.  Strong approximation gives density
of $S(k)$ in $J_T$, hence normality of $C$.  It also gives density of
$\Gamma$ in $K^T$.  By \eqref{arith:lattice-index}, the closure of
$N\cap\Gamma$ has finite index in $K^T$ and is therefore open.  Thus $C$
is open in $J_T$.  For every $v\notin T$, the group
$C\cap S(k_v)$ is an open normal subgroup of $S(k_v)$.  It is noncentral,
and $S$ is isotropic at $v$ because $A\subset T$.  Local simplicity and
perfection imply $S(k_v)\subset C$.  Finite-support products are dense
in the restricted product, so $C=J_T$.

Every coset of $N$ in $S(k)$ is now dense in $J_T$ and hence meets the
open subgroup $K^T$.  Thus every such coset meets $\Gamma$, and
\[
 [S(k):N]\leq[\Gamma:N\cap\Gamma]<\infty,
\]
as required.
\end{proof}

For an integer $e\geq1$, let
\begin{equation}\label{arith:power-notation}
 \begin{split}
 R&=S(k)^e:=\langle g^e:g\in S(k)\rangle,\qquad
 P_0=\overline{\delta_A(R)}\subset H_A,\\
 V_0&=\delta_A^{-1}(P_0),\qquad V=V_0/R.
 \end{split}
\end{equation}
The notation $S(k)^e$ means the subgroup generated by powers, not merely
the image of the power map.  The distinction will be essential when the
characteristic divides $e$.

\begin{proposition}\label{arith:power-defect}
The subgroup $R$ has finite index in $S(k)$; $P_0$ is open and normal in
$H_A$; and $V$ is finite.  Conjugation by $U(k)$ preserves $R$ and $V_0$,
and hence acts on $V$.  If $V=1$ for every integer $e\geq1$, then the
Margulis--Platonov assertion holds for $S(k)$.
\end{proposition}

\begin{proof}
The map $g\mapsto g^e$ is dominant: over an algebraic closure it is
surjective on each maximal torus, and conjugates of a maximal torus are
dense in $S$.  This dominance does not require separability of the power
map.  Weak approximation implies that $S(k)$ is Zariski dense, so the
set of rational $e$th powers is Zariski dense as well.  The characteristic
subgroup $R$ is therefore noncentral, and
\cref{arith:finite-index} gives its finite index.

Weak approximation gives density of $\delta_A(S(k))$ in $H_A$.  A finite
coset decomposition of $S(k)$ by $R$ then gives a finite coset covering
of $H_A$ by the closed subgroup $P_0$.  A closed subgroup of finite index
is open.  Normality follows from density and normality of $R$, and the
finiteness of $V$ follows from that of $S(k)/R$.
Finally, $U(k)$ normalizes $S(k)$ and preserves its power subgroup.
Taking closures at $A$ shows that it also preserves $P_0$ and $V_0$.
In fact weak approximation for $U$ shows that $P_0$ is normalized by
$\prod_{v\in A}U(k_v)$.

For the last assertion, let $N\triangleleft S(k)$ be noncentral.
By \cref{arith:finite-index}, the quotient $S(k)/N$ is finite.
Choose $e$ divisible by its exponent; then $R\subset N$.
If $V=1$, we have $R=\delta_A^{-1}(P_0)$, and density induces an
isomorphism $S(k)/R\simeq H_A/P_0$.  The inverse image in $H_A$ of
the subgroup corresponding to $N/R$ is open and normal and has rational
inverse image $N$.  This proves the claimed reduction.
\end{proof}

Our task is to prove $V=1$ for every $e$.  The quotient $V$ measures
precisely the failure of the power subgroup $R$ to contain all rational
points allowed by its closure at the anisotropic places.  The next
proposition provides the approximation in each individual coset of $R$;
passing only to its image in a finite quotient would not suffice.
The power-subgroup reduction and this use of exact cosets follow the
organization of \cite[Section~2.2]{MPNF}; the arithmetic inputs above
are their function-field forms.

\begin{proposition}[Closure of an exact coset]\label{arith:closure}
Let $v_0\notin A$, and write $S(\bbA^{v_0})$ for the restricted product
of $S(k_v)$ over $v\ne v_0$.  Then
\begin{equation}\label{arith:adelic-closure}
 \overline{R}^{\,S(\bbA^{v_0})}
 =\{g\in S(\bbA^{v_0}):g_A\in P_0\}.
\end{equation}
For any $\gamma\in S(k)$, the closure of $\gamma R$ is obtained by
replacing $P_0$ on the right by $\delta_A(\gamma)P_0$.  In particular,
all cosets $\gamma R$ with $\gamma\in V_0$ have the same closure.

Consequently one may approximate in an actual prescribed coset
$\gamma R$, meeting prescribed nonempty local open sets at finitely many
isotropic places and a prescribed nonempty open subset of
$\delta_A(\gamma)P_0$ at $A$.
If $x_*\in U(k)$, the same assertion holds after left translation to
the determinant slice $x_*S$ and to the coset $x_*\gamma R$.
In each assertion one may also require that the resulting point lie in
finitely many nonempty $k$-Zariski open subsets of $S$ or $x_*S$,
respectively.
\end{proposition}

\begin{proof}
Strong approximation gives density of $S(k)$ in $S(\bbA^{v_0})$.  Since
$R$ is a normal subgroup of finite index, its closure $C$ is a closed
normal subgroup of finite index, hence an open subgroup.  At every
isotropic place $v\ne v_0$, the intersection $C\cap S(k_v)$ is open,
normal, and noncentral.  The local result recalled above gives
$S(k_v)\subset C$.  Taking closures of finite-support products, $C$
contains the entire restricted product of these isotropic factors.

The remaining factors form the finite compact product $H_A$.
It follows that $C=P\times\prod_{v\notin A\cup\{v_0\}}'S(k_v)$ for a
closed subgroup $P\subset H_A$.  Since $R\subset C$, we have
$P_0\subset P$.  Conversely, the inverse image of $P_0$ is closed and
contains $R$, so it contains $C$ and gives $P\subset P_0$.  This proves
\eqref{arith:adelic-closure}.  Translation gives the assertion for
$\gamma R$, and left translation by $x_*$ gives the determinant-slice
version.  For finite local conditions that include $v_0$, apply the same
argument with a different isotropic place outside the specified finite
set omitted instead.  Finally, a nonempty Zariski open subset of a
smooth geometrically irreducible variety over a local field meets every
nonempty analytic open subset.  Impose the Zariski conditions at one
place, adding a place of approximation if necessary, and shrink its
local open set.  The same density argument then enforces them on the
rational point.
\end{proof}

\subsection{Finite quotients of groups already understood}

The induction will restrict the finite quotient $S(k)/R$ to smaller
special unitary groups, often over finite extensions of $k$.  The precise
consequence of their Margulis--Platonov theorem is the following.

\begin{proposition}\label{arith:finite-quotients}
Let $F$ be a global function field and $H$ a simply connected absolutely
almost simple $F$-group for which the Margulis--Platonov assertion holds.
Put $A_H=\{w:\rank_{F_w}H=0\}$.  Every homomorphism
$\psi:H(F)\to Q$ to a finite group extends uniquely to a continuous
homomorphism
\[
 \widehat\psi:\prod_{w\in A_H}H(F_w)\longrightarrow Q
\]
with the same image.  In particular, $\psi$ kills every rational point
whose coordinates at $A_H$ are sufficiently close to the identity.
If $A_H=\varnothing$, every such $\psi$ is trivial.
The same conclusions hold factor by factor for finite products of
restriction-of-scalars groups.
\end{proposition}

\begin{proof}
Strong approximation for $H$ away from an auxiliary split place gives
weak approximation by the same argument as in
\cref{arith:unitary-reduction}.  In particular, $H(F)$ is Zariski dense.
The kernel of $\psi$ has finite index, so connectedness of $H$ makes
its Zariski closure all of $H$ as well.  Thus this kernel is noncentral.  By the assumed theorem,
\[
 \ker\psi=\delta_{A_H}^{-1}(W)
\]
for an open normal subgroup $W$ of
$\prod_{w\in A_H}H(F_w)$.  This product is compact, so its quotient by
$W$ is finite.  Weak approximation gives density of $H(F)$ in the
product, and hence a canonical isomorphism
\[
 H(F)/\ker\psi\ \simeq\
 \left(\prod_{w\in A_H}H(F_w)\right)/W.
\]
Compose the quotient map with this isomorphism and the induced injection
into $Q$.  This constructs $\widehat\psi$; density proves uniqueness and
equality of images.  An identity neighborhood contained in $W$ gives
the local smallness assertion.
For a restriction of scalars, rational points and local points are the
corresponding groups over $F$ and products over the places above a given
place.  For a finite product, the images of the individual factors
commute.  Their continuous extensions therefore combine to give the
stated factorwise conclusion.
\end{proof}
